\section{The sample itself is a random object}

In descriptive statistics, we begin with observed values
\[
x_1,x_2,\ldots,x_n.
\]
These values are already known. We can arrange them, draw a histogram, compute
their mean, or describe their empirical distribution.

Before the observations are collected, however, there is no fixed list of
numbers. There is a random procedure that will produce such a list. We represent
its possible results by
\[
X_1,X_2,\ldots,X_n.
\]
The whole sample is therefore the random vector
\[
(X_1,X_2,\ldots,X_n).
\]
It is one random object with many possible values.

\begin{definitionbox}
A \textbf{random sample} is a collection of random observations
\[
X_1,X_2,\ldots,X_n
\]
generated according to a specified probabilistic model. Before observation, the
sample has many possible values. After observation, one of these possibilities
becomes the observed sample
\[
(x_1,x_2,\ldots,x_n).
\]
\end{definitionbox}

The word ``sample'' is therefore used at two related levels. Before the data are
known, it means a random object. After the data are known, it means one realized
list of values. The notation helps us keep these levels separate:
\[
(X_1,\ldots,X_n)
\quad\longrightarrow\quad
(x_1,\ldots,x_n).
\]
The arrow represents observation. It does not mean that randomness disappears
from the model. It means that one possible random sample has become the sample
we actually saw.

\subsection*{A small sample has a distribution}

Consider a mechanism that produces either failure, recorded as \(0\), or success,
recorded as \(1\). Suppose that the probability of success is
\[
P(X_i=1)=p,
\]
and the probability of failure is
\[
P(X_i=0)=1-p.
\]
Now take two observations and record their order. The random sample is
\[
(X_1,X_2).
\]
Its possible values are
\[
(0,0),\qquad (0,1),\qquad (1,0),\qquad (1,1).
\]

If the two observations are independent repetitions of the same mechanism, then
these possible samples have probabilities
\[
P((X_1,X_2)=(0,0))=(1-p)^2,
\]
\[
P((X_1,X_2)=(0,1))=(1-p)p,
\]
\[
P((X_1,X_2)=(1,0))=p(1-p),
\]
\[
P((X_1,X_2)=(1,1))=p^2.
\]
They add to one:
\[
(1-p)^2+(1-p)p+p(1-p)+p^2=1.
\]

This is a probability distribution on possible samples. The elementary outcomes
are no longer single zeros or ones. They are complete samples of size two.

\begin{examplebox}
Let \(p=0.6\). Then
\[
P((0,0))=0.16,\qquad P((0,1))=0.24,
\]
\[
P((1,0))=0.24,\qquad P((1,1))=0.36.
\]
Before sampling, all four samples are possible. After the experiment, we may
observe, for example,
\[
(x_1,x_2)=(1,0).
\]
The observed pair is one realization of the random pair \((X_1,X_2)\).
\end{examplebox}

The samples \((0,1)\) and \((1,0)\) contain the same number of successes, but they
are different ordered samples. They become indistinguishable only if we decide to
record a summary such as the total number of successes.

\subsection*{The sampling model determines the law of the sample}

The probabilities above did not come from the list of possible samples alone.
They came from assumptions about how the observations were generated. We assumed
that each observation had the same success probability \(p\) and that the two
observations were independent.

In introductory models, we often assume that
\[
X_1,X_2,\ldots,X_n
\]
are independent and have the same probability law. This is commonly abbreviated
by saying that the observations are \textbf{iid}: independent and identically
distributed.

These assumptions are useful, but they should not be treated as automatic. If
observations influence one another, or if objects are sampled without
replacement from a small population, then the possible samples or their
probabilities may change. The law of the sample always belongs to a specified
sampling mechanism.

\begin{remarkbox}
A dataset does not carry its sampling model inside itself. The model must state
how the observations were generated. The same observed values may arise from
different random mechanisms.
\end{remarkbox}

\subsection*{Three different distributions}

At this point it is important to separate three objects.

First, the \textbf{population law} describes one random observation, for example
\[
P(X_i=1)=p.
\]

Second, the \textbf{law of the sample} describes the random vector
\[
(X_1,\ldots,X_n).
\]
It assigns probabilities to possible complete samples.

Third, after one sample has been observed, its \textbf{empirical distribution}
describes the values that occurred in that particular dataset.

For example, if the observed sample is
\[
(1,0,1,1,0),
\]
then its empirical distribution assigns relative frequency \(3/5\) to \(1\) and
\(2/5\) to \(0\). This empirical distribution describes one realized sample. It
is not the probability law of all possible samples.

\begin{summarybox}
The sampling process creates a new probability space:
\[
\text{probability model for one observation}
\longrightarrow
\text{probability law for complete samples}
\longrightarrow
\text{one observed dataset}.
\]
A statistic will be obtained by applying a rule to each possible sample. Because
the sample is random, the statistic will also be random.
\end{summarybox}
